\documentclass[11pt,a4paper]{article}
\usepackage[utf8]{inputenc}
\usepackage[T1]{fontenc}
\usepackage{lmodern}
\usepackage{amsmath}
\usepackage[margin=2cm]{geometry}
\usepackage{enumitem}
\pagestyle{plain}
\setlength{\parindent}{0pt}
\setlength{\parskip}{0pt}
\begin{document}
{\Large\bfseries Tarea de peróxidos e hidróxidos}
\vspace{12pt}
\begin{enumerate}[label=\textbf{\arabic*.},leftmargin=*,itemsep=11pt,topsep=0pt]
\item \begin{minipage}[t]{\linewidth}
¿Cuál es la fórmula del peróxido de magnesio?

\vspace{5pt}
\begin{tabular}{@{}p{0.47\linewidth}p{0.47\linewidth}@{}}
a)~$\mathrm{Mg_{2}O_{2}}$ & b)~$\mathrm{MgO_{2}}$ \\
c)~$\mathrm{Mg_{2}(O_{2})_{2}}$ & d)~MgO \\
\end{tabular}
\end{minipage}

\item \begin{minipage}[t]{\linewidth}
¿Cómo se llama el $\mathrm{BaO_{2}}$?

\vspace{5pt}
\begin{tabular}{@{}p{0.47\linewidth}p{0.47\linewidth}@{}}
a)~peróxido de bario & b)~óxido bárico \\
c)~óxido de bario (IV) & d)~hidróxido de bario \\
\end{tabular}
\end{minipage}

\item \begin{minipage}[t]{\linewidth}
¿Cuál es el nombre sistemático del $\mathrm{Fe_{2}(O_{2})_{3}}$?

\vspace{5pt}
\begin{tabular}{@{}p{0.47\linewidth}p{0.47\linewidth}@{}}
a)~trióxido de dihierro & b)~triperóxido de dihierro \\
c)~peróxido de dihierro & d)~diperóxido de trihierro \\
\end{tabular}
\end{minipage}

\item \begin{minipage}[t]{\linewidth}
¿Cuál es el nombre tradicional del $\mathrm{CuO_{2}}$?

\vspace{5pt}
\begin{tabular}{@{}p{0.47\linewidth}p{0.47\linewidth}@{}}
a)~óxido cúprico & b)~óxido cuproso \\
c)~peróxido cuproso & d)~peróxido cúprico \\
\end{tabular}
\end{minipage}

\item \begin{minipage}[t]{\linewidth}
¿Cuál es el nombre stock del $\mathrm{NiCl_{2}}$?

\vspace{5pt}
\begin{tabular}{@{}p{0.47\linewidth}p{0.47\linewidth}@{}}
a)~dicloruro de níquel (II) & b)~clorato de níquel (II) \\
c)~cloruro de níquel (III) & d)~cloruro de níquel (II) \\
\end{tabular}
\end{minipage}

\item \begin{minipage}[t]{\linewidth}
¿Cuál es la fórmula del bromuro de magnesio?

\vspace{5pt}
\begin{tabular}{@{}p{0.47\linewidth}p{0.47\linewidth}@{}}
a)~$\mathrm{Mg_{2}Br}$ & b)~MgBr \\
c)~$\mathrm{MgBr_{2}}$ & d)~$\mathrm{Br_{2}Mg}$ \\
\end{tabular}
\end{minipage}

\item \begin{minipage}[t]{\linewidth}
¿Cuál es la fórmula del hidróxido cuproso?

\vspace{5pt}
\begin{tabular}{@{}p{0.47\linewidth}p{0.47\linewidth}@{}}
a)~$\mathrm{Cu(OH)_{2}}$ & b)~CuOH \\
c)~$\mathrm{Cu(OH)_{3}}$ & d)~$\mathrm{Cu_{2}OH}$ \\
\end{tabular}
\end{minipage}

\item \begin{minipage}[t]{\linewidth}
¿Cuál es el nombre sistemático del $\mathrm{Sn(OH)_{4}}$?

\vspace{5pt}
\begin{tabular}{@{}p{0.47\linewidth}p{0.47\linewidth}@{}}
a)~tetrahidróxido de estaño & b)~tetraóxido de estaño \\
c)~hidróxido estannoso & d)~hidróxido de estaño (II) \\
\end{tabular}
\end{minipage}

\item \begin{minipage}[t]{\linewidth}
¿Cuál es el nombre tradicional del $\mathrm{Al(OH)_{3}}$?

\vspace{5pt}
\begin{tabular}{@{}p{0.47\linewidth}p{0.47\linewidth}@{}}
a)~hidróxido alumínico & b)~trihidróxido alumínico \\
c)~hidróxido de aluminio (III) & d)~hidróxido aluminoso \\
\end{tabular}
\end{minipage}

\item \begin{minipage}[t]{\linewidth}
¿Cuál es el nombre stock del $\mathrm{Au_{2}O_{2}}$?

\vspace{5pt}
\begin{tabular}{@{}p{0.47\linewidth}p{0.47\linewidth}@{}}
a)~óxido de oro (I) & b)~óxido de oro (II) \\
c)~peróxido de oro (I) & d)~peróxido de oro (II) \\
\end{tabular}
\end{minipage}

\end{enumerate}
\end{document}
